\documentclass[10pt]{extarticle}
\usepackage[a4paper,margin=1.5cm]{geometry}
\usepackage[T1]{fontenc}
\usepackage{amsmath,amssymb}
\usepackage{graphicx}
\usepackage[export]{adjustbox}
\usepackage{multicol}
\usepackage{enumitem}
\usepackage{booktabs,array,tabularx}
\setlength{\columnsep}{0.7cm}
\setlength{\parindent}{0pt}
\setlength{\parskip}{3pt}
\pagestyle{empty}
\makeatletter
% figure registry (figdefs.tex is generated by the extract step)
\newcommand{\qfigdef}[4]{\@namedef{qf@file@#1}{#2}\@namedef{qf@w@#1}{#3}\@namedef{qf@wide@#1}{#4}}
\newcommand{\qf@check}[1]{\@ifundefined{qf@file@#1}{\PackageError{score2}{Unknown figure id #1}{}}{}}
\newcommand{\qf@img}[4][]{\includegraphics[width=#2\dimexpr\@nameuse{qf@w@#4}pt\relax,max width=#3,max height=0.45\textheight,#1]{\@nameuse{qf@file@#4}}}
% Full-width questions: finalize wraps every question that contains a wide figure in \qfullbegin...\qfullend,
% so the WHOLE question (stem, figures, options) is set across the page in reading order, then the two columns
% resume. The question's head - its number, stem and first figure (\qfullsplit marks the end of the head) -
% is kept on one page, so a number is never left at a page foot with its figure overleaf; the rest may break
% across pages normally, so no big gaps are left. Without a \qfullsplit, the whole question is kept together
% if it is under 45% of a page. Never write these yourself.
\newif\ifqf@full
\newif\ifqf@boxed
\newsavebox{\qf@box}
\newcommand{\qf@open}{\global\qf@fulltrue\par\global\qf@boxedtrue
  \setbox\qf@box\vbox\bgroup\hsize\textwidth\linewidth\textwidth}
\newcommand{\qf@emit}[1]{\par\egroup\global\qf@boxedfalse
  \ifdim\dimexpr\ht\qf@box+\dp\qf@box\relax<#1\textheight\noindent\usebox{\qf@box}\par\else\unvbox\qf@box\fi}
\newcommand{\qfullsplit}{\ifqf@boxed\qf@emit{0.85}\fi}
\newcommand{\qf@close}{\ifqf@boxed\qf@emit{0.45}\fi\global\qf@fullfalse}
\newcommand{\qfullbegin}{\end{multicols}\qf@open}
\newcommand{\qfullend}{\qf@close\begin{multicols}{2}}
\newcommand{\qfullstart}{\qf@open}
\newcommand{\qfulldone}{\qf@close}
% \qfig[scale]{id}: block figure, centred, directly where it is written (size set by the tool).
% Fallback if a wide figure is not inside a full-width question: break out of the columns in place - no floats.
\newcommand{\qfig}[2][1]{\qf@check{#2}%
  \ifqf@full\qf@blk{#1}{#2}\else
  \ifnum\@nameuse{qf@wide@#2}=1 \expandafter\@firstoftwo\else\expandafter\@secondoftwo\fi
  {\end{multicols}\begin{center}\qf@img{#1}{\textwidth}{#2}\end{center}\begin{multicols}{2}}%
  {\qf@blk{#1}{#2}}\fi}
\newcommand{\qf@blk}[2]{\par\smallskip{\centering\qf@img{#1}{\linewidth}{#2}\par}\smallskip}
% \qfigi[scale]{id}: inline figure (inside option lists / tables / minipages); never breaks columns.
\newcommand{\qfigi}[2][1]{\qf@check{#2}\qf@img[valign=c]{#1}{\linewidth}{#2}}
\makeatother
% structure macros - lint reads these to check completeness and figure placement
\newcommand{\papersec}[1]{\par\bigskip{\centering\Large\bfseries PAPER-#1\par}\medskip}
\newcommand{\partsec}[2]{\par\medskip\noindent{\large\bfseries PART-#1: #2}\par\smallskip}
\newcommand{\sect}[1]{\par\smallskip\noindent\textbf{#1}\par}
\newcommand{\qn}[1]{\par\smallskip\noindent\textbf{#1.}~}
\newcommand{\paraq}[2]{\par\smallskip\noindent\textbf{Paragraph for Questions #1 and #2}\par}
\newenvironment{qoptions}{\begin{itemize}[leftmargin=1.8em,itemsep=1pt,topsep=2pt,label={}]}{\end{itemize}}
\newcommand{\opt}[2]{\item[(#1)] #2}
% figure registry (generated by score2_tools.py finalize)
\qfigdef{p1_phy_s1_q01_1}{Score Advanced - Test 26 (21 Apr 2024)-img/p1_phy_s1_q01_1.png}{151.6}{0}
\qfigdef{p1_phy_s1_q02_1}{Score Advanced - Test 26 (21 Apr 2024)-img/p1_phy_s1_q02_1.png}{74.9}{0}
\qfigdef{p1_phy_s1_q03_1}{Score Advanced - Test 26 (21 Apr 2024)-img/p1_phy_s1_q03_1.png}{59.9}{0}
\qfigdef{p1_phy_s1_q04_1}{Score Advanced - Test 26 (21 Apr 2024)-img/p1_phy_s1_q04_1.png}{113.7}{0}
\qfigdef{p1_phy_s2_q05_1}{Score Advanced - Test 26 (21 Apr 2024)-img/p1_phy_s2_q05_1.png}{97.9}{0}
\qfigdef{p1_phy_s2_q06_1}{Score Advanced - Test 26 (21 Apr 2024)-img/p1_phy_s2_q06_1.png}{59.9}{0}
\qfigdef{p1_phy_s2_q07_1}{Score Advanced - Test 26 (21 Apr 2024)-img/p1_phy_s2_q07_1.png}{56.6}{0}
\qfigdef{p1_phy_s3_q08_1}{Score Advanced - Test 26 (21 Apr 2024)-img/p1_phy_s3_q08_1.png}{63.0}{0}
\qfigdef{p1_phy_s3_q08_2}{Score Advanced - Test 26 (21 Apr 2024)-img/p1_phy_s3_q08_2.png}{66.8}{0}
\qfigdef{p1_phy_s3_q08_3}{Score Advanced - Test 26 (21 Apr 2024)-img/p1_phy_s3_q08_3.png}{68.5}{0}
\qfigdef{p1_phy_s3_q08_4}{Score Advanced - Test 26 (21 Apr 2024)-img/p1_phy_s3_q08_4.png}{105.6}{0}
\qfigdef{p1_phy_s3_q10_1}{Score Advanced - Test 26 (21 Apr 2024)-img/p1_phy_s3_q10_1.png}{108.2}{0}
\qfigdef{p1_phy_s3_q11_1}{Score Advanced - Test 26 (21 Apr 2024)-img/p1_phy_s3_q11_1.png}{134.9}{0}
\qfigdef{p1_phy_s4_q02_1}{Score Advanced - Test 26 (21 Apr 2024)-img/p1_phy_s4_q02_1.png}{42.3}{0}
\qfigdef{p1_phy_s4_q03_1}{Score Advanced - Test 26 (21 Apr 2024)-img/p1_phy_s4_q03_1.png}{97.5}{0}
\qfigdef{p1_phy_s4_q04_1}{Score Advanced - Test 26 (21 Apr 2024)-img/p1_phy_s4_q04_1.png}{192.7}{0}
\qfigdef{p1_phy_s4_q05_1}{Score Advanced - Test 26 (21 Apr 2024)-img/p1_phy_s4_q05_1.png}{137.4}{0}
\qfigdef{p1_chem_s1_q01_A}{Score Advanced - Test 26 (21 Apr 2024)-img/p1_chem_s1_q01_A.png}{122.0}{0}
\qfigdef{p1_chem_s1_q01_B}{Score Advanced - Test 26 (21 Apr 2024)-img/p1_chem_s1_q01_B.png}{121.7}{0}
\qfigdef{p1_chem_s1_q01_C}{Score Advanced - Test 26 (21 Apr 2024)-img/p1_chem_s1_q01_C.png}{122.0}{0}
\qfigdef{p1_chem_s1_q01_D}{Score Advanced - Test 26 (21 Apr 2024)-img/p1_chem_s1_q01_D.png}{121.7}{0}
\qfigdef{p1_chem_s1_q04_A}{Score Advanced - Test 26 (21 Apr 2024)-img/p1_chem_s1_q04_A.png}{132.5}{0}
\qfigdef{p1_chem_s1_q04_B}{Score Advanced - Test 26 (21 Apr 2024)-img/p1_chem_s1_q04_B.png}{142.5}{0}
\qfigdef{p1_chem_s1_q04_C}{Score Advanced - Test 26 (21 Apr 2024)-img/p1_chem_s1_q04_C.png}{181.5}{0}
\qfigdef{p1_chem_s1_q04_D}{Score Advanced - Test 26 (21 Apr 2024)-img/p1_chem_s1_q04_D.png}{182.0}{0}
\qfigdef{p1_chem_s2_q06_B}{Score Advanced - Test 26 (21 Apr 2024)-img/p1_chem_s2_q06_B.png}{135.7}{0}
\qfigdef{p1_chem_s2_q06_C}{Score Advanced - Test 26 (21 Apr 2024)-img/p1_chem_s2_q06_C.png}{168.9}{0}
\qfigdef{p1_chem_s2_q06_D}{Score Advanced - Test 26 (21 Apr 2024)-img/p1_chem_s2_q06_D.png}{101.4}{0}
\qfigdef{p1_chem_s3_q09_1}{Score Advanced - Test 26 (21 Apr 2024)-img/p1_chem_s3_q09_1.png}{53.6}{0}
\qfigdef{p1_chem_s3_q10_1}{Score Advanced - Test 26 (21 Apr 2024)-img/p1_chem_s3_q10_1.png}{83.2}{0}
\qfigdef{p1_chem_s3_q10_2}{Score Advanced - Test 26 (21 Apr 2024)-img/p1_chem_s3_q10_2.png}{85.8}{0}
\qfigdef{p1_chem_s3_q10_3}{Score Advanced - Test 26 (21 Apr 2024)-img/p1_chem_s3_q10_3.png}{83.2}{0}
\qfigdef{p1_chem_s3_q10_4}{Score Advanced - Test 26 (21 Apr 2024)-img/p1_chem_s3_q10_4.png}{55.1}{0}
\qfigdef{p1_chem_s3_q10_5}{Score Advanced - Test 26 (21 Apr 2024)-img/p1_chem_s3_q10_5.png}{55.1}{0}
\qfigdef{p1_chem_s4_q03_1}{Score Advanced - Test 26 (21 Apr 2024)-img/p1_chem_s4_q03_1.png}{230.7}{0}
\qfigdef{p1_chem_s4_q03_2}{Score Advanced - Test 26 (21 Apr 2024)-img/p1_chem_s4_q03_2.png}{225.8}{0}
\qfigdef{p1_chem_s4_q04_1}{Score Advanced - Test 26 (21 Apr 2024)-img/p1_chem_s4_q04_1.png}{232.4}{0}
\qfigdef{p2_phy_s1_q01_1}{Score Advanced - Test 26 (21 Apr 2024)-img/p2_phy_s1_q01_1.png}{103.1}{0}
\qfigdef{p2_phy_s1_q04_1}{Score Advanced - Test 26 (21 Apr 2024)-img/p2_phy_s1_q04_1.png}{169.0}{0}
\qfigdef{p2_phy_s2_q05_1}{Score Advanced - Test 26 (21 Apr 2024)-img/p2_phy_s2_q05_1.png}{161.3}{0}
\qfigdef{p2_phy_s2_q06_1}{Score Advanced - Test 26 (21 Apr 2024)-img/p2_phy_s2_q06_1.png}{103.6}{0}
\qfigdef{p2_phy_s2_q07_1}{Score Advanced - Test 26 (21 Apr 2024)-img/p2_phy_s2_q07_1.png}{97.7}{0}
\qfigdef{p2_phy_s3_para3_4_1}{Score Advanced - Test 26 (21 Apr 2024)-img/p2_phy_s3_para3_4_1.png}{108.5}{0}
\qfigdef{p2_phy_s4_q05_1}{Score Advanced - Test 26 (21 Apr 2024)-img/p2_phy_s4_q05_1.png}{92.1}{0}
\qfigdef{p2_phy_s4_q06_1}{Score Advanced - Test 26 (21 Apr 2024)-img/p2_phy_s4_q06_1.png}{97.2}{0}
\qfigdef{p2_phy_s4_q08_1}{Score Advanced - Test 26 (21 Apr 2024)-img/p2_phy_s4_q08_1.png}{83.3}{0}
\qfigdef{p2_phy_s4_q09_1}{Score Advanced - Test 26 (21 Apr 2024)-img/p2_phy_s4_q09_1.png}{76.8}{0}
\qfigdef{p2_phy_s4_q10_1}{Score Advanced - Test 26 (21 Apr 2024)-img/p2_phy_s4_q10_1.png}{184.7}{0}
\qfigdef{p2_chem_s1_q01_A}{Score Advanced - Test 26 (21 Apr 2024)-img/p2_chem_s1_q01_A.png}{74.9}{0}
\qfigdef{p2_chem_s1_q01_B}{Score Advanced - Test 26 (21 Apr 2024)-img/p2_chem_s1_q01_B.png}{74.0}{0}
\qfigdef{p2_chem_s1_q01_C}{Score Advanced - Test 26 (21 Apr 2024)-img/p2_chem_s1_q01_C.png}{74.4}{0}
\qfigdef{p2_chem_s1_q01_D}{Score Advanced - Test 26 (21 Apr 2024)-img/p2_chem_s1_q01_D.png}{74.0}{0}
\qfigdef{p2_chem_s1_q03_A}{Score Advanced - Test 26 (21 Apr 2024)-img/p2_chem_s1_q03_A.png}{137.0}{0}
\qfigdef{p2_chem_s1_q03_B}{Score Advanced - Test 26 (21 Apr 2024)-img/p2_chem_s1_q03_B.png}{143.2}{0}
\qfigdef{p2_chem_s1_q03_C}{Score Advanced - Test 26 (21 Apr 2024)-img/p2_chem_s1_q03_C.png}{126.8}{0}
\qfigdef{p2_chem_s1_q03_D}{Score Advanced - Test 26 (21 Apr 2024)-img/p2_chem_s1_q03_D.png}{145.8}{0}
\qfigdef{p2_chem_s2_q06_1}{Score Advanced - Test 26 (21 Apr 2024)-img/p2_chem_s2_q06_1.png}{77.4}{0}
\qfigdef{p2_chem_s3_para3_4_F}{Score Advanced - Test 26 (21 Apr 2024)-img/p2_chem_s3_para3_4_F.png}{286.6}{1}
\qfigdef{p2_chem_s4_q05_1}{Score Advanced - Test 26 (21 Apr 2024)-img/p2_chem_s4_q05_1.png}{134.5}{0}
\qfigdef{p2_chem_s4_q08_1}{Score Advanced - Test 26 (21 Apr 2024)-img/p2_chem_s4_q08_1.png}{111.9}{0}
\qfigdef{p2_chem_s4_q09_1}{Score Advanced - Test 26 (21 Apr 2024)-img/p2_chem_s4_q09_1.png}{192.3}{0}
\raggedcolumns
\makeatletter\@beginparpenalty=10000\makeatother
\newcolumntype{L}[1]{>{\raggedright\arraybackslash}p{#1}}
\begin{document}
\begin{center}{\Large\bfseries Score Advanced --- Test 26 (21 Apr 2024)}\end{center}
\begin{multicols}{2}
\raggedcolumns
\papersec{1}
\partsec{1}{PHYSICS}
\sect{SECTION-I (i)}
\qn{1} A nonconducting hollow cone has charge density 12$\varepsilon_0$ C/m$^2$. A part ABP is cut and removed from the cone. The radius of the base of the cone is 1m. The potential due to the remaining portion of the cone at point `P' is\par
\qfig{p1_phy_s1_q01_1}
\begin{qoptions}
\opt{A}{5V}
\opt{B}{6V}
\opt{C}{8V}
\opt{D}{10V}
\end{qoptions}
\qn{2} A square of side length a and mass 2m is cut across its diagonal. Two ideal light rods of negligible length are rigidly attached to the two corners that were cut. A tangential impulse J (into the page) is applied to the outermost corner of the square. What is the tension developed in each rod?\par
\qfig{p1_phy_s1_q02_1}
\begin{qoptions}
\opt{A}{$\dfrac{9(2-\sqrt2)J^2}{ma}$}
\opt{B}{$\dfrac{27(2-\sqrt2)J^2}{ma}$}
\opt{C}{$\dfrac{3J^2}{\sqrt2ma}$}
\opt{D}{$\dfrac{3J^2}{2\sqrt2ma}$}
\end{qoptions}
\qn{3} A meteor hits the moon and creates a crater within the ground. The crater can be modelled as a cubic hole of length L, created when the material originally within it is ejected out after the collision. Assume that energy is not lost during the collision, and kinetic energy of meteor gets converted to gravitational potential energy of the ejected mass. If a meteor of twice the mass of the initial meteor hits the moon at the same velocity, how much larger is the volume of the crater compared to the initial one? Assume everything on the moon are of the same density.\par
\qfig{p1_phy_s1_q03_1}
\begin{qoptions}
\opt{A}{16}
\opt{B}{4}
\opt{C}{2}
\opt{D}{$2^{3/4}$}
\end{qoptions}
\qn{4} On one of the plates of the capacitors has a small hole of radius r, covered with a soap film. The surface tension of the film is $\sigma$. The capacitor is charged to a potential V. Distance d between the plates is small compared to linear dimensions of the plates. Value of h ($\ll$ r) is equal to :-\par
\qfig{p1_phy_s1_q04_1}
\begin{qoptions}
\opt{A}{$\dfrac{\varepsilon_0V^2r^2}{8\sigma d^2}$}
\opt{B}{$\dfrac{\varepsilon_0V^2r^2}{4\sigma d^2}$}
\opt{C}{$\dfrac{\varepsilon_0V^2r^2}{16\sigma d^2}$}
\opt{D}{$\dfrac{\varepsilon_0V^2r^2}{\sigma d^2}$}
\end{qoptions}
\sect{SECTION-I (ii)}
\qn{5} The length of a weightless rod is 2$\ell$. A small sphere of mass m is fixed at a distance $\ell$ from its upper end. The rod rests with its one end against the wall and the other against the floor. The end that rests on the floor is being moved with a constant velocity V away from the wall :-\par
\qfig{p1_phy_s2_q05_1}
\begin{qoptions}
\opt{A}{The force with which the sphere affects the rod at the moment when $\alpha=45^\circ$ is m$\left(g-\dfrac{V^2}{\sqrt2\ell}\right)$.}
\opt{B}{The force with which the sphere affects the rod at the moment when $\alpha=45^\circ$ is m$\left(g-\dfrac{V^2}{2\sqrt2\ell}\right)$.}
\opt{C}{Vertical acceleration of m when $\alpha=60^\circ$ is $\dfrac{2V^2}{\ell}$.}
\opt{D}{Horizontal acceleration of m when $\alpha=60^\circ$ is $\dfrac{V^2}{\ell}$.}
\end{qoptions}
\qn{6} Three adjacent walls of a regular hexagon-shaped room are covered with flat mirror (see the figure). A lamp is lit at the centre of the room.\par
\qfig{p1_phy_s2_q06_1}
\begin{qoptions}
\opt{A}{The number of images seen by an observer at A is 4.}
\opt{B}{The number of images seen by an observer at B is 4.}
\opt{C}{The number of images seen by an observer at C is 5.}
\opt{D}{The number of images seen by an observer at D is 5.}
\end{qoptions}
\qn{7} At the end of the thermally insulated cylinder with the piston there is a valve (figure) that covers a small hole, which opens and begins to let air from outside into the cylinder at a pressure difference $\Delta P=P_0/3$ ($P_0$ is atmospheric pressure). The valve does not allow air to escape from the cylinder to the outside. At the initial moment of time, the piston is pressed against the end of the cylinder, there is no air inside.\par
In the first case, the cylinder is filled with air to volume $V_0$, slowly moving the piston, after which it is stopped, and then the piston is released.\par
In the second case, the piston is abruptly moved to a position in which the volume under the piston is equal to $V_0$, so that the air does not have time to penetrate through the valve into the cylinder. In this position, fix the piston, wait until the cylinder is filled with air, just as in the first case, release the piston. Out side the cylinder is air at atmospheric pressure $P_0$ and temperature $T_0$. The friction of the piston against the walls, the mass of the piston, as well as the heat exchange of air with the piston and the cylinder walls can be neglected. Air can be considered a diatomic ideal gas. After the piston is released, the valve remains closed all the time.\par
\qfig{p1_phy_s2_q07_1}
\begin{qoptions}
\opt{A}{For first case the temperature of air inside the cylinder at the moment the piston is stopped at volume $V_0$ is $T_0$.}
\opt{B}{For first case, the temperature of air inside the cylinder when the piston stops movement is $\dfrac87T_0$.}
\opt{C}{For second case, the temperature of the air in the cylinder after stopping of piston and filling with air is $\dfrac75T_0$.}
\opt{D}{For second case, the temperature of the air after piston is released and it stops movement is $\dfrac95T_0$.}
\end{qoptions}
\sect{SECTION-I (iii)}
\qn{8} List$-$I contains four different YDSE systems and List$-$II contains intensity measured at some points on the screen. Do the correct match (es) in List$-$I and List$-$II, (CP$_1$ = 0.3 mm and CP$_2$ = 1.2 mm)\par
{\footnotesize\setlength{\tabcolsep}{1.5pt}\noindent\begin{tabular}{|c|L{5.1cm}|c|L{1.4cm}|} \hline & \multicolumn{1}{c|}{List$-$I} & & \multicolumn{1}{c|}{List$-$II}\\ \hline
(P) & \qfigi[0.75]{p1_phy_s3_q08_1}\par $\lambda$ = 4000 \AA\ in air\par d = 1mm, D = 1m\par Intensity due to each slit = $I_0$ & (1) & $I_C=4I_0$\\ \hline
(Q) & \qfigi[0.75]{p1_phy_s3_q08_2}\par $\lambda$ = 4000 \AA\ in air\par d = 1mm, D = 1m\par Intensity due to each slit = $I_0$ and no absorption of light by water ($\mu_W$ = 4/3) & (2) & $I_{P_1}=2I_0$\\ \hline
(R) & \qfigi[0.75]{p1_phy_s3_q08_3}\par $\lambda$ = 4000 \AA\ in air\par d = 1mm, D = 1m\par Film thickness t = 0.8 $\mu$m, refractive index of film = 3/2\par Intensity due to each slit = $I_0$ and no absorption of light by glass & (3) & $I_{P_1}=4I_0$\\ \hline
(S) & \qfigi[0.75]{p1_phy_s3_q08_4}\par $\lambda$ = 4000 \AA\ in air\par d = 1mm, D = 1m\par Intensity due to each slit = $I_0$ & (4) & $I_{P_1}=0$\\ \hline
 & & (5) & $I_{P_2}=2I_0$\\ \hline
\end{tabular}}\par
\begin{qoptions}
\opt{A}{P $\rightarrow$ 1,2;Q $\rightarrow$ 1,3;R $\rightarrow$ 2;S $\rightarrow$ 4,5}
\opt{B}{P $\rightarrow$ 1,2;Q $\rightarrow$ 1,3;R $\rightarrow$ 2,5;S $\rightarrow$ 4,5}
\opt{C}{P $\rightarrow$ 2;Q $\rightarrow$ 3;R $\rightarrow$ 2,5;S $\rightarrow$ 1,2}
\opt{D}{P $\rightarrow$ 1,2;Q $\rightarrow$ 1,3;R $\rightarrow$ 1,2;S $\rightarrow$ 4,5}
\end{qoptions}
\qn{9} An electron in a hydrogen atom makes a transition $n_1\rightarrow n_2$, where $n_1$ and $n_2$ are the principal quantum numbers of the two states. Assume Bohr model to be valid.\par
{\footnotesize\setlength{\tabcolsep}{1.5pt}\noindent\begin{tabular}{|c|L{3.9cm}|c|L{1.9cm}|} \hline & \multicolumn{1}{c|}{List$-$I} & & \multicolumn{1}{c|}{List$-$II}\\ \hline
(P) & The electron emits an energy of 2.55 eV & (1) & $n_1=2,n_2=1$\\ \hline
(Q) & Time period of the electron in the initial state is eight times that in the final state. & (2) & $n_1=4,n_2=2$\\ \hline
(R) & Speed of electron become two times & (3) & $n_1=5,n_2=3$\\ \hline
(S) & Radius of orbit of electron becomes 4.77 \AA & (4) & $n_1=6,n_2=3$\\ \hline
 & & (5) & $n_1=8,n_2=4$\\ \hline
\end{tabular}}\par
\begin{qoptions}
\opt{A}{P $\rightarrow$ 2;Q $\rightarrow$ 1,2;R $\rightarrow$ 1,2;S $\rightarrow$ 3}
\opt{B}{P $\rightarrow$ 2;Q $\rightarrow$ 1,2,4;R $\rightarrow$ 1,3,4;S $\rightarrow$ 4}
\opt{C}{P $\rightarrow$ 2;Q $\rightarrow$ 1,2,4,5;R $\rightarrow$ 1,2,4,5;S $\rightarrow$ 3,4}
\opt{D}{P $\rightarrow$ 2;Q $\rightarrow$ 1,2,3,4,5;R $\rightarrow$ 1,2,3,4,5;S $\rightarrow$ 4}
\end{qoptions}
\qn{10} A load is supported using three wires of same cross section area as shown.\par
\qfig{p1_phy_s3_q10_1}
{\footnotesize\setlength{\tabcolsep}{1.5pt}\noindent\begin{tabular}{|c|L{3.9cm}|c|L{1.9cm}|} \hline & \multicolumn{1}{c|}{List$-$I} & & \multicolumn{1}{c|}{List$-$II}\\ \hline
(P) & Equal tensile stress in all wires & (1) & $Y_2=2Y_1$\\ \hline
(Q) & Equal tension in all wires & (2) & $Y_2=4Y_1$\\ \hline
(R) & Equal elastic potential energy in all wires & (3) & $Y_2=8Y_1$\\ \hline
(S) & Equal elastic potential energy per unit volume in all wires & (4) & $Y_2=16Y_1$\\ \hline
 & & (5) & Does not depend on ratio of $\dfrac{Y_1}{Y_2}$\\ \hline
\end{tabular}}\par
\begin{qoptions}
\opt{A}{P $\rightarrow$ 2;Q $\rightarrow$ 2;R $\rightarrow$ 3;S $\rightarrow$ 4}
\opt{B}{P $\rightarrow$ 2;Q $\rightarrow$ 3;R $\rightarrow$ 2;S $\rightarrow$ 4}
\opt{C}{P $\rightarrow$ 4;Q $\rightarrow$ 5;R $\rightarrow$ 3;S $\rightarrow$ 1}
\opt{D}{P $\rightarrow$ 1;Q $\rightarrow$ 2;R $\rightarrow$ 5;S $\rightarrow$ 4}
\end{qoptions}
\qn{11} V$-$T graph of a thermodynamic process is shown in figure. If heat given to process is taken to be positive :-\par
($\Delta$Q is heat given to the gas in the process, $\Delta$U is change in internal energy in the process, $\Delta$W is the work done by the gas in the process).\par
\qfig{p1_phy_s3_q11_1}
{\footnotesize\setlength{\tabcolsep}{1.5pt}\noindent\begin{tabular}{|c|L{2.0cm}|c|L{2.0cm}|} \hline & \multicolumn{1}{c|}{List$-$I} & & \multicolumn{1}{c|}{List$-$II}\\ \hline
(P) & JK & (1) & $\Delta W\geq0$\\ \hline
(Q) & KL & (2) & $\Delta Q<0$\\ \hline
(R) & LM & (3) & $\Delta W<0$\\ \hline
(S) & MJ & (4) & $\Delta Q>0$\\ \hline
 & & (5) & $\Delta U=0$\\ \hline
\end{tabular}}\par
\begin{qoptions}
\opt{A}{P $\rightarrow$ 2,3,5;Q $\rightarrow$ 4;R $\rightarrow$ 1,4,5;S $\rightarrow$ 2,3}
\opt{B}{P $\rightarrow$ 1,3,5;Q $\rightarrow$ 4;R $\rightarrow$ 1,4,5;S $\rightarrow$ 2,5}
\opt{C}{P $\rightarrow$ 2,3,5;Q $\rightarrow$ 1,4;R $\rightarrow$ 1,4,5;S $\rightarrow$ 2,3}
\opt{D}{P $\rightarrow$ 1,4;Q $\rightarrow$ 2,4;R $\rightarrow$ 1,3,5;S $\rightarrow$ 2,3}
\end{qoptions}
\sect{SECTION-II}
\qn{1} The pitch of a screw gauge is 1 mm and its cap is divided into 100 divisions. When nothing is placed between its studs, the zero of the circular scale lies 6 divisions below the reference line and zero of main scale is just visible. When a wire is placed between its studs, the main scale reading is 4 mm and 56th division of circular scale coincides with the reference line. The length of the wire is 4.20 cm. If the curved surface area (in cm$^2$) comes out to be x, then find the value of 100x. $\left(\pi=\dfrac{22}{7}\right)$
\qn{2} A tower is glued together from three identical wooden blocks of edge lengths a and b as in the figure below. What is the largest ratio a/b such that the tower would still be stable?\par
\qfig{p1_phy_s4_q02_1}
\qn{3} I placed 2 identical cylinder-shaped barrels, having the same height h, on the horizontal ground, one on the top of the other (see the figure). Both were nearly fully filled with water. Your task is to drill a hole perpendicularly to the wall of each of the barrels at the same xh height measured from the bottom of each barrel. The value of the dimensionless factor x, in order that the impact points of the water rays fall as far apart as possible. The value of $\dfrac1x$ is :-\par
\qfig{p1_phy_s4_q03_1}
\qn{4} Two infinitely long current carrying wires carry constant current $i_1$ and $i_2$ as shown in the diagram. The equations of the wire curvatures are $y^2-8x-6y+25=0$ and x = 0. The magnitude of force acting on one of the wires due to the other is $\dfrac{\mu_0i_1i_2}{\beta}$. The value of $\beta$ is (The current-carrying wires are rigidly fixed. The units for distances on the graph should be taken in meters.)\par
\qfig{p1_phy_s4_q04_1}
\qn{5} Consider two convex lens with focal length $f_a$ and $f_b$ respectively. If placed appropriately, a small object perpendicular to the optic axis can create an inverted image, the size and position of which matches the size and position of object as shown in the figure. The distance between the lens is d and distance between the object and the left lens is u. The ratio $\dfrac ud$ is (Take : $f_a$ = 1 cm, $f_b$ = 4 cm)\par
\qfig{p1_phy_s4_q05_1}
\qn{6} A capacitor of value 1.0 F discharges through a device whose resistance R varies linearly with applied potential difference, V, so that R = AV + B, where A and B are constants. The resistance of the device has a value of 10.0$\Omega$ when V = 6.0 V, and 4.0$\Omega$ when V = 0.06 V. The capacitor is initially charged to a potential of 6.0V. Determine how long (in sec) takes for the capacitor to discharge to 1\% of this initial value. (Take : $\ell$n10 = 2)
\partsec{2}{CHEMISTRY}
\sect{SECTION-I (i)}
\qn{1} Select the correct option for following aqueous solutions :\par
\begin{qoptions}
\opt{A}{\qfigi{p1_chem_s1_q01_A}}
\opt{B}{\qfigi{p1_chem_s1_q01_B}}
\opt{C}{\qfigi{p1_chem_s1_q01_C}}
\opt{D}{\qfigi{p1_chem_s1_q01_D}}
\end{qoptions}
\qn{2} H$_3$C $-$ C $\equiv$ C $-$ Ph $\xrightarrow{\text{Hg}^{2+}/\text{H}^+/\text{H}_2\text{O}}$ P, P is :
\begin{qoptions}
\opt{A}{H$_3$C $-$ C(=O) $-$ CH$_2-$Ph}
\opt{B}{H$_3$C $-$ CH$_2-$ C(=O) $-$ Ph}
\opt{C}{H$_3$C $-$ C(=O) $-$ C(=O) $-$ Ph}
\opt{D}{H$_3$C $-$ CH(OH) $-$ CH$_2-$Ph}
\end{qoptions}
\qn{3} Alkaline solution of (X) $\xrightarrow{\text{NaBO}_3}$ Canary yellow colored solution\par
Acidic solution of (Y) $\xrightarrow{\text{NaBiO}_3}$ Lilac colored solution\par
(X) \& (Y) are respectively :-
\begin{qoptions}
\opt{A}{Fe(OH)$_3$ , Mn(OH)$_2$}
\opt{B}{Cr(OH)$_3$ , Pb(OH)$_2$}
\opt{C}{Fe(OH)$_3$ , Zn(OH)$_2$}
\opt{D}{Cr(OH)$_3$ , Mn(OH)$_2$}
\end{qoptions}
\qn{4} Among the following, which reaction will not give N-ethylcyclopentylamine as major product?
\begin{qoptions}
\opt{A}{\qfigi{p1_chem_s1_q04_A}}
\opt{B}{\qfigi{p1_chem_s1_q04_B}}
\opt{C}{\qfigi{p1_chem_s1_q04_C}}
\opt{D}{\qfigi{p1_chem_s1_q04_D}}
\end{qoptions}
\sect{SECTION-I (ii)}
\qn{5} Consider two liquids A and B having pure vapour pressures $P_A^\circ$ \& $P_B^\circ$ forming an ideal solution. The plot of $\dfrac1{X_A}$ v/s $\dfrac1{Y_A}$ (where $X_A$ and $Y_A$ are the mole fraction of liquid A in liquid and vapour phase respectively) is linear with slope and Y intercepts respectively :
\begin{qoptions}
\opt{A}{Intercept = $\dfrac{P_A^\circ-P_B^\circ}{P_B^\circ}$}
\opt{B}{Intercept = $\dfrac{P_B^\circ-P_A^\circ}{P_B^\circ}$}
\opt{C}{Slope = $\dfrac{P_A^\circ}{P_B^\circ}$}
\opt{D}{Slope = $\dfrac{P_B^\circ}{P_A^\circ}$}
\end{qoptions}
\qn{6} In which of the following reactions cinnamic acid is one of the major organic product ?\par
\begin{qoptions}
\opt{A}{Ph$-$CH$_2-$COOEt $\xrightarrow[\text{(iii)}\Delta]{\text{(i)EtONa; (ii)H}_3\text{O}^\oplus}$}
\opt{B}{\qfigi{p1_chem_s2_q06_B}}
\opt{C}{\qfigi{p1_chem_s2_q06_C}}
\opt{D}{\qfigi{p1_chem_s2_q06_D}}
\end{qoptions}
\qn{7} The complex(es) with distorted octahedron structure is/are :-
\begin{qoptions}
\opt{A}{[Cu(H$_2$O)$_6$]$^{2+}$}
\opt{B}{[FeF$_6$]$^{3-}$}
\opt{C}{[MnF$_6$]$^{3-}$}
\opt{D}{[Fe(CN)$_6$]$^{4-}$}
\end{qoptions}
\sect{SECTION-I (iii)}
\qn{8} Four different gases and their Vander Waal's constant are given below : (R = 0.08 L-atm/K-mol)\par
{\footnotesize\setlength{\tabcolsep}{3pt}\noindent\begin{tabular}{|c|c|c|} \hline Gas & a (in atm L$^2$ / mol$^2$) & b (in L/mol)\\ \hline
W & 1.5 & 0.05\\ \hline
X & 2.0 & 0.08\\ \hline
Y & 4.2 & 0.025\\ \hline
Z & 5.0 & 0.10\\ \hline
\end{tabular}}\par
{\footnotesize\setlength{\tabcolsep}{1.5pt}\noindent\begin{tabular}{|c|c|c|L{5.6cm}|} \hline \multicolumn{2}{|c|}{List - I} & \multicolumn{2}{c|}{List - II}\\ \hline
(P) & W & (1) & Gas having highest intermolecular force of attraction.\\ \hline
(Q) & X & (2) & Gas having smallest size of molecules.\\ \hline
(R) & Y & (3) & Gas that can be liquefied at 85 K\\ \hline
(S) & Z & (4) & Gas resembles with ideal gas at nearly boiling point of water\\ \hline
 & & (5) & Gas with highest boyle's temperature.\\ \hline
\end{tabular}}\par
Codes :
\begin{qoptions}
\opt{A}{P $\rightarrow$ 3,4,5;Q $\rightarrow$ 1,5;R $\rightarrow$ 2,4;S $\rightarrow$ 5}
\opt{B}{P $\rightarrow$ 1,2;Q $\rightarrow$ 4,5;R $\rightarrow$ 3,4;S $\rightarrow$ 2}
\opt{C}{P $\rightarrow$ 4;Q $\rightarrow$ 3;R $\rightarrow$ 2,5;S $\rightarrow$ 1}
\opt{D}{P $\rightarrow$ 1,4;Q $\rightarrow$ 3,4;R $\rightarrow$ 5;S $\rightarrow$ 1,2,3}
\end{qoptions}
\qn{9} For the reaction, using suitable alkyl halide and a suitable nucleophile what will be the correct match for the nucleophile used as per List-I and the solvent being used in List-II?\par
{\footnotesize\setlength{\tabcolsep}{1.5pt}\noindent\begin{tabular}{|c|L{2.3cm}|c|L{3.6cm}|} \hline & \multicolumn{1}{c|}{List-I} & & \multicolumn{1}{c|}{List-II}\\ \hline
(P) & MeS$^-$ & (1) & Reacts with CH$_3$CH$_2$Cl to give S$_\text{N}$2 product\\ \hline
(Q) & NC$^-$ + NH$_4$Cl & (2) & Reacts with CH$_3$CHO followed by hydrolysis to give alanine.\\ \hline
(R) & \qfigi[0.75]{p1_chem_s3_q09_1} & (3) & Reacts with EtX followed by NH$_2$NH$_2$ gives EtNH$_2$\\ \hline
(S) & Ph$_3$P: + nBuLi & (4) & Reacts with CH$_3$CH$_2$Cl followed by acetone to give 2-methyl-but-2-ene.\\ \hline
 & & (5) & Reacts with chloro benzene followed by H$_3$O$^+$ give aniline.\\ \hline
\end{tabular}}\par
\begin{qoptions}
\opt{A}{P $\rightarrow$ 1;Q $\rightarrow$ 2;R $\rightarrow$ 3;S $\rightarrow$ 4}
\opt{B}{P $\rightarrow$ 5;Q $\rightarrow$ 2;R $\rightarrow$ 1;S $\rightarrow$ 4}
\opt{C}{P $\rightarrow$ 5;Q $\rightarrow$ 4;R $\rightarrow$ 3;S $\rightarrow$ 1}
\opt{D}{P $\rightarrow$ 1;Q $\rightarrow$ 4;R $\rightarrow$ 5;S $\rightarrow$ 3}
\end{qoptions}
\qn{10} {\footnotesize\setlength{\tabcolsep}{1.5pt}\noindent\begin{tabular}{|c|L{3.5cm}|c|L{2.0cm}|} \hline & \multicolumn{1}{c|}{\parbox{3.5cm}{\centering List$-$I (Molecules/ions)}} & & \multicolumn{1}{c|}{\parbox{2.0cm}{\centering List$-$II (Arrangement of given entity in List-I)}}\\ \hline
(P) & CO on each manganese in [Mn$_2$(CO)$_{10}$] & (1) & Triangular pyramidal\\ \hline
(Q) & Cl on each Al in Al$_2$Cl$_6$ & (2) & Octahedral\\ \hline
(R) & Two S around each S in S$_8$ molecule & (3) & Square pyramidal\\ \hline
(S) & Oxygen atoms around each sulphur in S$_2$O$_6^{2-}$ & (4) & Tetrahedral\\ \hline
 & & (5) & Bent\\ \hline
\end{tabular}}\par
Select \textbf{CORRECT} code :-
\begin{qoptions}
\opt{A}{P $\rightarrow$ 2;Q $\rightarrow$ 4;R $\rightarrow$ 5;S $\rightarrow$ 1}
\opt{B}{P $\rightarrow$ 3;Q $\rightarrow$ 4;R $\rightarrow$ 5;S $\rightarrow$ 1}
\opt{C}{P $\rightarrow$ 3;Q $\rightarrow$ 1;R $\rightarrow$ 2;S $\rightarrow$ 4}
\opt{D}{P $\rightarrow$ 2;Q $\rightarrow$ 2;R $\rightarrow$ 4;S $\rightarrow$ 5}
\end{qoptions}
\qn{11} {\footnotesize\setlength{\tabcolsep}{1.5pt}\noindent\begin{tabular}{|c|L{2.6cm}|c|c|} \hline & \multicolumn{1}{c|}{\parbox{2.6cm}{\centering List-I (Element in periodic table)}} & & \multicolumn{1}{c|}{\parbox{2.6cm}{\centering List-II [Graphical change in radius (r) with atomic number (Z)]}}\\ \hline
(P) & Halogens & (1) & \qfigi[0.75]{p1_chem_s3_q10_1}\\ \hline
(Q) & M$^{+3}$ ions of Lanthanides & (2) & \qfigi[0.75]{p1_chem_s3_q10_2}\\ \hline
(R) & M atoms of Lanthanides & (3) & \qfigi[0.75]{p1_chem_s3_q10_3}\\ \hline
(S) & Size of 3d series metal (covalent radius) & (4) & \qfigi[0.75]{p1_chem_s3_q10_4}\\ \hline
 & & (5) & \qfigi[0.75]{p1_chem_s3_q10_5}\\ \hline
\end{tabular}}\par
Select \textbf{CORRECT} code :-
\begin{qoptions}
\opt{A}{P $\rightarrow$ 4;Q $\rightarrow$ 5;R $\rightarrow$ 1;S $\rightarrow$ 2}
\opt{B}{P $\rightarrow$ 5;Q $\rightarrow$ 1;R $\rightarrow$ 3;S $\rightarrow$ 2}
\opt{C}{P $\rightarrow$ 4;Q $\rightarrow$ 3;R $\rightarrow$ 1;S $\rightarrow$ 2}
\opt{D}{P $\rightarrow$ 5;Q $\rightarrow$ 1;R $\rightarrow$ 2;S $\rightarrow$ 3}
\end{qoptions}
\sect{SECTION-II}
\qn{1} Identify correct number of statements :\par
\textbf{Statement - 1 :} Entropy of system increases in dissociation of ammonia gas into N$_2$ and H$_2$ gases.\par
\textbf{Statement - 2 :} Entropy of system increases in polymerization process of ethene to form polythene.\par
\textbf{Statement - 3 :} Entropy of system increases in adiabatic irreversible compression of an ideal gas.\par
\textbf{Statement - 4 :} Entropy of system increases in condensation of H$_2$O (vapour).\par
\textbf{Statement - 5 :} Entropy of system increases in isothermal free expansion of ideal gas.\par
\textbf{Statement - 6 :} Entropy of system always increases in a spontaneous process.\par
\textbf{Statement - 7 :} Adiabatic process whether reversible or irreversible will not cause any change in entropy of the surroundings.\par
\textbf{Statement - 8 :} At absolute zero, entropy of any substance becomes zero.
\qn{2} 1.2 gram of carbon is burnt completely with 2.56 gram of oxygen to produce CO(g) and CO$_2$(g). None of the reactant remaining after reaction. Produced mixture of gases is treated with solid I$_2$O$_5$ (to know the amount of CO(g) produced). The liberated iodine required x mL of 0.1 M hypo solution for complete titration. The value of x is \_\_\_\_\_\_\_\_
\qn{3} An organic compound reacts with excess MeMgBr followed by H$_2$O to give two organic product of which one reacts with NaOI to give yellow precipitate and OTHER DOES NOT. How many of the following compounds can satisfy this condition ?\par
\qfig{p1_chem_s4_q03_1}
\qfig{p1_chem_s4_q03_2}
\qn{4} Consider the following reactions and write molecular weight of \textbf{white precipitate (Y)}.\par
(Given atomic mass : H = 1, O = 16, C = 12, F = 19, N = 14, S = 32, P = 31, Cl = 35.5, Br = 80, I = 127, K = 39, Na = 23, Pb = 207, Al = 27, Si = 28, Mn = 55, Cr = 52)\par
\qfig{p1_chem_s4_q04_1}
Colored gas [Q] is identified by its staining of starch paper orange red\par
Phenol + [Q] in water $\rightarrow$ \textbf{White precipitate (Y)} + other product
\qn{5} Find the Mo\textemdash Mo bond order in [($\eta^5$-C$_5$H$_5$)Mo(CO)$_2$]$_2$\par
(Given that both the Mo have 18 valence electrons in given complex)\par
[Atomic number of Mo = 42]
\qn{6} Consider all possible isomeric ketones including stereoisomers of MW = 86, All these isomers are independently reacted with NaBH$_4$ (NOTE : stereoisomers are also reacted separately).\par
The total number of ketones that give a racemic product(s) is/are.
\partsec{3}{MATHEMATICS}
\sect{SECTION-I (i)}
\qn{1} If $I=\displaystyle\int_{1/4}^{4}\dfrac{\tan^{-1}x}{x}dx$, then the value of 2[I] is equal to - (where [.] denotes greatest integer function and $e^{0.693}\simeq2$)
\begin{qoptions}
\opt{A}{6}
\opt{B}{4}
\opt{C}{2}
\opt{D}{8}
\end{qoptions}
\qn{2} If $\displaystyle\int_0^1xf(x^2)\,dx=\dfrac12$, then $\displaystyle\int_1^2f^{-1}(x)\,dx$ is equal to (where f(0) = 1, f(1) = 2 \& f is strictly increasing in its domain)
\begin{qoptions}
\opt{A}{0}
\opt{B}{2}
\opt{C}{1}
\opt{D}{3}
\end{qoptions}
\qn{3} Number of solutions of $\sqrt{\tan\theta}=2\sin\theta$, $\theta\in[0,2\pi]$ is equal to
\begin{qoptions}
\opt{A}{2}
\opt{B}{4}
\opt{C}{5}
\opt{D}{7}
\end{qoptions}
\qn{4} Value of $\displaystyle\lim_{n\to\infty}n\int_{0^+}^{\pi/2}\left(1-\sqrt[n]{\sin x}\right)dx$ is
\begin{qoptions}
\opt{A}{$\ell$n 2}
\opt{B}{$-\ell$n 2}
\opt{C}{$\dfrac\pi2\ell$n 2}
\opt{D}{$-\dfrac\pi2\ell$n 2}
\end{qoptions}
\sect{SECTION-I (ii)}
\qn{5} If $A_i$ is the area bounded by $|x-a_i|+|y|=b_i$, $i\in N$, where $a_{i+1}=a_i+\dfrac32b_i$ \& $b_{i+1}=\dfrac{b_i}{2}$, $a_1=0$, $b_1=32$, then
\begin{qoptions}
\opt{A}{$A_3=128$}
\opt{B}{$A_3=256$}
\opt{C}{$\displaystyle\lim_{n\to\infty}\sum_{i=1}^nA_i=\dfrac83(32)^2$}
\opt{D}{$\displaystyle\lim_{n\to\infty}\sum_{i=1}^nA_i=\dfrac43(16)^2$}
\end{qoptions}
\qn{6} A circle S having centre ($\alpha$, $\beta$) intersect a parabola $y^2=4x$ at three points A, B and C such that normals at these points are concurrent at (9, 6), where O is origin. The slope of normals at A \& B are positive, then
\begin{qoptions}
\opt{A}{Sum of modulus of slopes of normals at points A, B and C is 6}
\opt{B}{$|\alpha-\beta|=4$}
\opt{C}{point of intersection of tangents at A \& B will be (2, $-$3)}
\opt{D}{Circle S also passes through O}
\end{qoptions}
\qn{7} If $P=\displaystyle\sum_{r=3n}^{4n-1}\left[\dfrac{r^2+13n^2-7rn}{n^3}\right]$ \& $Q=\displaystyle\sum_{r=3n+1}^{4n}\left[\dfrac{r^2+13n^2-7rn}{n^3}\right]$, then
\begin{qoptions}
\opt{A}{$P>\dfrac56$}
\opt{B}{$Q<\dfrac56$}
\opt{C}{$Q>\dfrac56$}
\opt{D}{$P<\dfrac56$}
\end{qoptions}
\sect{SECTION-I (iii)}
\qn{8} Match List-I with List-II and select the correct answer using the code given below the list.\par
{\footnotesize\setlength{\tabcolsep}{1.5pt}\noindent\begin{tabular}{|c|L{5.0cm}|c|L{0.5cm}|} \hline & \multicolumn{1}{c|}{List-I} & & \multicolumn{1}{c|}{List-II}\\ \hline
(P) & Given all terms of an A.P. are positive integers, third term is 8 and sum of first 15 terms lies between 191 and 196, then first term is & (1) & 1\\ \hline
(Q) & Suppose that for some angles x and y the equation $\cos2y-\cos2x=3a-2$ and $\cos^2x+\sin^2y=\dfrac{a^2}{2}$ hold simultaneously, then number of possible values of a is/are & (2) & 4\\ \hline
(R) & Number of solution(s) of the equation $\sin^{-1}\dfrac{2x}{1+x^2}=\dfrac\pi2+\cos^{-1}x$ & (3) & 6\\ \hline
(S) & In $\Delta$ABC, minimum value of $\dfrac1\Delta\sqrt{abc\,(a+b+c)}$ is, (all symbols have their standard meaning) & (4) & 0\\ \hline
 & & (5) & 5\\ \hline
\end{tabular}}\par
\begin{qoptions}
\opt{A}{P $\rightarrow$ 1;Q $\rightarrow$ 1;R $\rightarrow$ 1;S $\rightarrow$ 2}
\opt{B}{P $\rightarrow$ 3;Q $\rightarrow$ 5;R $\rightarrow$ 2;S $\rightarrow$ 3}
\opt{C}{P $\rightarrow$ 3;Q $\rightarrow$ 1;R $\rightarrow$ 2;S $\rightarrow$ 1}
\opt{D}{P $\rightarrow$ 3;Q $\rightarrow$ 1;R $\rightarrow$ 1;S $\rightarrow$ 2}
\end{qoptions}
\qn{9} Match List-I with List-II and select the correct answer using the code given below the list.\par
{\footnotesize\setlength{\tabcolsep}{1.5pt}\noindent\begin{tabular}{|c|L{5.0cm}|c|L{0.5cm}|} \hline & \multicolumn{1}{c|}{List-I} & & \multicolumn{1}{c|}{List-II}\\ \hline
(P) & Given $\displaystyle\int_{-2}^2f(\{x\})dx=0$ and $\displaystyle\int_{-2}^6f(\{x\})dx=-2$ then value of $\displaystyle\int_6^2f(\{x\})dx$ is equal to (Here \{.\} denotes fractional part of function) & (1) & 2\\ \hline
(Q) & If $|\vec a-\vec b|=8$, $|\vec b|=\sqrt{57}$ and $|\vec a+\vec b|=10$, then value of $|\vec a|$ is equal to & (2) & 3\\ \hline
(R) & $\displaystyle\int_{-\pi}^{\pi}\dfrac{|\text{sgn}(\cot x)+\text{sgn}(\tan x)|\,dx}{\pi}$ is & (3) & 4\\ \hline
(S) & Let $A=\begin{bmatrix}1&\frac12\\0&1\end{bmatrix}\begin{bmatrix}1&\frac1{2^2}\\0&1\end{bmatrix}\begin{bmatrix}1&\frac1{2^3}\\0&1\end{bmatrix}\ldots\infty$ times, then sum of all elements of matrix A is & (4) & 5\\ \hline
 & & (5) & 6\\ \hline
\end{tabular}}\par
\begin{qoptions}
\opt{A}{P $\rightarrow$ 5;Q $\rightarrow$ 4;R $\rightarrow$ 3;S $\rightarrow$ 2}
\opt{B}{P $\rightarrow$ 1;Q $\rightarrow$ 4;R $\rightarrow$ 3;S $\rightarrow$ 2}
\opt{C}{P $\rightarrow$ 1;Q $\rightarrow$ 4;R $\rightarrow$ 2;S $\rightarrow$ 2}
\opt{D}{P $\rightarrow$ 1;Q $\rightarrow$ 2;R $\rightarrow$ 3;S $\rightarrow$ 2}
\end{qoptions}
\qn{10} Suppose n people (n $>$ 6) are asked a question successively in a random order and exactly 3 out of `n' people know the answer. Let $P_r$ denote the probability that $r^{\text{th}}$ person asked, is the first to know the answer, then\par
{\footnotesize\setlength{\tabcolsep}{1.5pt}\noindent\begin{tabular}{|c|L{2.6cm}|c|L{3.2cm}|} \hline & \multicolumn{1}{c|}{List-I} & & \multicolumn{1}{c|}{List-II}\\ \hline
(P) & Probability that first four do not know the answer & (1) & $\dfrac{3(n-3)}{n(n-1)}$\\ \hline
(Q) & $P_2$ & (2) & $\dfrac{(n-4)(n-5)(n-6)}{n(n-1)(n-2)}$\\ \hline
(R) & $P_{n-2}$ & (3) & $\dfrac{6}{n(n-1)(n-2)}$\\ \hline
(S) & $P_r$ & (4) & $\dfrac{3(n-r)(n-r-1)}{n(n-1)(n-2)}$\\ \hline
 & & (5) & $\dfrac{6(n-r)(n-r-1)}{n(n-1)(n-2)}$\\ \hline
\end{tabular}}\par
\begin{qoptions}
\opt{A}{P $\rightarrow$ 2;Q $\rightarrow$ 1;R $\rightarrow$ 3;S $\rightarrow$ 4}
\opt{B}{P $\rightarrow$ 1;Q $\rightarrow$ 2;R $\rightarrow$ 3;S $\rightarrow$ 4}
\opt{C}{P $\rightarrow$ 2;Q $\rightarrow$ 1;R $\rightarrow$ 3;S $\rightarrow$ 5}
\opt{D}{P $\rightarrow$ 1;Q $\rightarrow$ 2;R $\rightarrow$ 4;S $\rightarrow$ 3}
\end{qoptions}
\qn{11} {\footnotesize\setlength{\tabcolsep}{1.5pt}\noindent\begin{tabular}{|c|L{4.5cm}|c|L{0.5cm}|} \hline & \multicolumn{1}{c|}{List-I} & & \multicolumn{1}{c|}{List-II}\\ \hline
(P) & $\displaystyle\lim_{x\to0}\left[\left(\min\left(y^2+2y+4\right)\right)\dfrac{\sin x}{x}\right]$ is equal to, where [ . ] denotes G.I.F. & (1) & 2\\ \hline
(Q) & $\displaystyle\lim_{x\to0}\dfrac1x\sin^{-1}\left(\dfrac{2x}{1+x^2}\right)=$ & (2) & 3\\ \hline
(R) & The number of points in the domain of the function $f(x)=\dfrac{1}{\ln|x-4|}$ is discontinuous are & (3) & 6\\ \hline
(S) & Number of critical points of f(x) = min (sinx, cosx), in x $\in$ [0, 2$\pi$], is & (4) & 0\\ \hline
 & & (5) & 5\\ \hline
\end{tabular}}\par
\begin{qoptions}
\opt{A}{P $\rightarrow$ 1;Q $\rightarrow$ 1;R $\rightarrow$ 2;S $\rightarrow$ 5}
\opt{B}{P $\rightarrow$ 1;Q $\rightarrow$ 1;R $\rightarrow$ 4;S $\rightarrow$ 5}
\opt{C}{P $\rightarrow$ 2;Q $\rightarrow$ 4;R $\rightarrow$ 5;S $\rightarrow$ 2}
\opt{D}{P $\rightarrow$ 2;Q $\rightarrow$ 4;R $\rightarrow$ 1;S $\rightarrow$ 3}
\end{qoptions}
\sect{SECTION-II}
\qn{1} If 4-horses participate in race. The number of ways in which they can go through finish line is $\alpha$ (Dead heat is possible), then value of $\alpha$ is (Two or more horses can finish race exactly at same time)
\qn{2} A line 3x + y = 8 touches a hyperbola H = 0 at P(1, 5) meets its asymptotes at A and B. If AB = $2\sqrt{10}$, C(1, 1) be the centre of hyperbola, e and $\ell$ are eccentricity and latus rectum, (here length of conjugate axis $<$ length of transverse axis), then value of [$e^2+\ell^2$] is (Where [.] Greatest Integer Function)
\qn{3} Lines $L_1:\dfrac{x-6}{3}=\dfrac{y-4}{2}=\dfrac{z-2}{0}$ and $L_2:\dfrac{x-8}{4}=\dfrac{y-2}{0}=\dfrac{z-4}{2}$ meets the plane $\pi:\vec r\cdot(2\hat i+\hat j-\hat k)=6$ at point A and B respectively. Area of the triangle formed by the lines $L_1$, $L_2$ and AB is
\qn{4} Let A and B are non-singular matrices of order 3 such that det. (A) = 5 and $A^{-1}B^2+AB=O$. If value of det.($A^6-2A^4B+A^2B^2$) is equal to $10^n$, then value of n is
\qn{5} Let $\omega$ be the complex number representing the point M $\left(\dfrac{-1}{2},\dfrac{\sqrt3}{2}\right)$, then number of distinct complex numbers z satisfying the equation $(z+1)\begin{vmatrix}z+\omega^2&1\\1&z+\omega\end{vmatrix}+\omega\begin{vmatrix}1&\omega\\z+\omega&\omega^2\end{vmatrix}+\omega^2\begin{vmatrix}\omega&z+\omega^2\\\omega^2&1\end{vmatrix}=0$ is equal to :
\qn{6} Let f(x) be a twice differentiable function defined on ($-\infty,\infty$) such that f(x) = f(2 $-$ x). Then value of $\displaystyle\int_{-1}^1f'(1+x)x^2e^{x^2}dx$ is equal to
\raggedcolumns
\papersec{2}
\partsec{1}{PHYSICS}
\sect{SECTION-I (i)}
\qn{1} Three particles A, B and C situated at vertices of an equilateral triangle, all moving with same constant speed such that A always move towards B, B always towards C and C always towards A. Initial separation between each of the particle is a. O is the centroid of the triangle. Distance covered by particle A when it completes one revolution around O is\par
\qfig{p2_phy_s1_q01_1}
\begin{qoptions}
\opt{A}{$2a\left(1-e^{-2\sqrt3\pi}\right)$}
\opt{B}{$\dfrac{2a}{3}\left(1-e^{-2\sqrt3\pi}\right)$}
\opt{C}{$a\left(1+e^{-2\sqrt3\pi}\right)$}
\opt{D}{$\dfrac{2a}{3}\left(1-e^{-\sqrt3\pi}\right)$}
\end{qoptions}
\qn{2} Use a diameter 1 mm superconductive wire to construct a circle of radius 5 cm. When the circular loop is superconductive, the current is 100 A. After one year, it is observed that the circular loop's change in current is less than $10^{-6}$ A. The resistivity of wire (in $\Omega$-m) is of the order of :-
\begin{qoptions}
\opt{A}{$10^{-20}$}
\opt{B}{$10^{-25}$}
\opt{C}{$10^{-30}$}
\opt{D}{$10^{-35}$}
\end{qoptions}
\qn{3} In cloudy weather, measurements are taken with a light meter pointed towards the sky. We find that, due to the scattering of light by the clouds, the incident power per unit area is approximately $I_0$, regardless of the orientation of the light meter. On top of an opaque thin walled spherical shell of radius R, there is a small hole of radius r (much larger than the wavelength of visible light). The inner side of the wall of the spherical shell is coated with soot. The intensity of illumination on the inner surface of the shell, when it is placed in the open air is :-
\begin{qoptions}
\opt{A}{$I_0\dfrac{r^2}{R^2}$}
\opt{B}{$I_0\dfrac{r^2}{4R^2}$}
\opt{C}{$I_0\dfrac{r^2}{R^2-r^2}$}
\opt{D}{$I_0\dfrac{r^2}{R^2+r^2}$}
\end{qoptions}
\qn{4} A massless thread makes N turns around statically fixed cylinder, as shown in the figure. Initially, the free (unwound) ends of the thread are parallel to the axis X. Then, a heavy point-like object P is attached to one end of the thread while the other end is pulled with a constant velocity u along X. Find the velocity attained by the heavy object when it detaches from the cylinder. Suppose that the turns of the thread are wound tightly to one another and are placed practically in the same plane, perpendicular to the cylinder axis. Neglect any friction in the system. Do not consider the force of gravity.\par
\qfig{p2_phy_s1_q04_1}
\begin{qoptions}
\opt{A}{u}
\opt{B}{$2\pi Nu$}
\opt{C}{$\sqrt{1+(2\pi N)^2}\,u$}
\opt{D}{$\pi Nu$}
\end{qoptions}
\sect{SECTION-I (ii)}
\qn{5} A car is moving with a velocity $\dfrac{3v}{10}$ towards a medium MNRP as shown, where `v' is the speed of sound in air. Two men A and B are moving towards each other in the medium with the velocities $\dfrac v6$ and $\dfrac v9$ respectively. The speed of sound in medium is 7v and medium is kept at rest. The car is emitting a sound of frequency f. Then choose the correct option(s).\par
\qfig{p2_phy_s2_q05_1}
\begin{qoptions}
\opt{A}{The wavelength received by the man A is $\dfrac{49v}{10f}$}
\opt{B}{The wavelength received by the man A is different from the wavelength received by man B}
\opt{C}{The ratio of frequency observed by the man A to that observed by the man B is $\dfrac{123}{128}$}
\opt{D}{The frequency observed by the man B is $\left(\dfrac{640}{441}\right)$f}
\end{qoptions}
\qn{6} A rod of mass `M' and length `L' was at rest in a vertical position, it's end A being hinged about which it can freely rotate. A slight disturbance caused the rod to start rotating clockwise. A bullet of mass `m' and speed `$v_0$' flying horizontally at a height `L/2' above ground hits the end B of the rod and gets embedded. (Take M = 3m). Then choose the correct statement(s).\par
\qfig{p2_phy_s2_q06_1}
\begin{qoptions}
\opt{A}{Minimum velocity of the bullet such that the rod again becomes vertical is $v_0=\sqrt{gL}\left(\sqrt6+\sqrt{20}\right)$}
\opt{B}{Tangential acceleration of the midpoint of the rod when it makes an angle $60^\circ$ with horizontal after collision is $\dfrac{5g}{16}$}
\opt{C}{Minimum velocity of the bullet such that the rod again becomes vertical is $v_0=\sqrt{3gL}\left(\sqrt6+\sqrt{20}\right)$}
\opt{D}{Tangential acceleration of the midpoint of the rod when it makes an angle $60^\circ$ with vertical after collision is $\dfrac{5g}{12}$}
\end{qoptions}
\qn{7} The equivalent circuit of a real source of alternating voltage of frequency $\omega=1.00\times10^3$ s$^{-1}$ consists of an ideal voltage source with the amplitude U = 15.0 V, a resistor with the resistance r = 10$\Omega$ and a capacitor with the capacitance C = 100 $\mu$F. Different circuits of resistors, capacitors and coils can be connected to the source as a load. A load circuit, which assures maximum heat generated in the load itself is connected across the source terminals AB.\par
\qfig{p2_phy_s2_q07_1}
\begin{qoptions}
\opt{A}{The load is purely resistive in nature.}
\opt{B}{The power factor of the load is $\dfrac1{\sqrt2}$.}
\opt{C}{The resistance of the load is equal to 10$\Omega$.}
\opt{D}{The maximum heat generated in the load is $\dfrac{45}{32}$W.}
\end{qoptions}
\sect{SECTION-II (i)}
\paraq{1}{2}
A solid cylinder of mass m and radius r rests on the inside of a thin-walled cylinder of mass M and radius R. The solid cylinder is made to oscillate around the equilibrium position.\par
(Take M = 1kg, m = 1 kg, R = 1m and r = 0.25 m)
\qn{1} The angular frequency (in rad/s) of oscillation if the large cylinder is fixed and small cylinder rolls without slipping in the bottom of the cylinder is :-
\qn{2} The angular frequency (in rad/s) of oscillation if the large cylinder is free to rotate about an axis through its fixed centre and small cylinder rolls without slipping in the bottom of the cylinder.
\paraq{3}{4}
The figure shows a diagram of an electrical circuit consisting of a source with electromotive force $\varepsilon$, a capacitor with capacitance C, a coil of inductance L, a key K and a metal disk of radius R. The disk can rotate around its axis without friction. The disk is located in a uniform magnetic field with induction B, directed along its axis as shown in the figure. Almost the entire mass of the disk is concentrated in its thin rim. The source is connected to the disk at its centre O, and the inductor is connected to the disk using a fixed sliding contact A. Initially, key K is open, the disk is motionless, there is no current in the coil, and capacitor is not charged. Key K is closed.\par
(Take : B = 1T, C = 1F, $\varepsilon$ = 1V, m = 1 kg, R = 1m, L = 1H, $\pi$ = 3.14 and $\sqrt5$ = 2.24)\par
\qfig{p2_phy_s3_para3_4_1}
\qn{3} The maximum angular velocity (in rad/sec) of the disc after key is closed is :-
\qn{4} The time (in sec) after closing the key after which the angular velocity of the disk first reaches the maximum value is :-
\sect{SECTION-II (ii)}
\qn{5} A ring of uniform mass density of mass m (m = 100 gm) and radius R (R = 50 cm) is suspended by two insulated strings PS and PT from the point P on y-axis in vertical plane. The ring is kept parallel to horizontal surface on x-z plane. Ring carries a current $\left(i=\dfrac7{22}\text{amp}\right)$ as shown in the figure. There exist a magnetic field (B = 1 Tesla) along negative x-axis. The tension in the string PS is $T_1$ and in the string PT is $T_2$. Find the ratio $\dfrac{T_1}{T_2}$. (It is given that $\theta=60^\circ$)\par
\qfig{p2_phy_s4_q05_1}
\qn{6} If we consider earth as a uniform solid sphere of mass M and Radius R then force of interaction between two hemi-spheres will be $\dfrac{3GM^2}{\alpha R^2}$. Find $\alpha$\par
\qfig{p2_phy_s4_q06_1}
\qn{7} A given sample contains two types of atoms A and B in the ratio 3 : 1. Atoms of type A undergo $\alpha$-decay with a half life of 30 days to form `B' while atoms of type B undergo $\alpha$-decay with a half life of 45 days to form `C', which is stable. Calculate the time (in days) after which the activities of A and that of B are in the ratio 9 : 22. If time is T write $\dfrac T{10}$.\par
$A\xrightarrow{T_{1/2}=30\text{ days}}B\xrightarrow{T_{1/2}=45\text{ days}}C$
\qn{8} Circular base of hemi-spherical shell base closed body is maintained at $T_o$ = 30K, outer atmosphere is at 0 K. If equilibrium temperature of curved surface is 10 $\times$ (x)$^{\frac34}$K, write x :-\par
\qfig{p2_phy_s4_q08_1}
\qn{9} The resistance between two joints is R. The equivalent resistance between A and B is $\dfrac{\alpha R}{15}$. The value of $\alpha$ is :-\par
\qfig{p2_phy_s4_q09_1}
\qn{10} Unpolarised light is incident on a water surface at the Brewster angle. The refracted ray meets the block of glass $\left[\mu_g=\dfrac32\right]$ immersed in the water $\left(\mu_w=\dfrac43\right)$. The flat upper surface of the block makes an angle $\alpha$ with water surface. The value of $\alpha$ for which the light reflected from the block is completely polarised. $\left(\tan^{-1}\dfrac98=48.4\right)$ Fill the value of 10$\alpha$.\par
\qfig{p2_phy_s4_q10_1}
\raggedcolumns
\partsec{2}{CHEMISTRY}
\sect{SECTION-I (i)}
\qn{1} The standard Gibb's free energy of the following reaction, $\Delta G^\circ$ have been determined at various temperatures :\par
(a) C(s) + O$_2$(g) $\rightarrow$ CO$_2$(g)\par
(b) C(s) + $\dfrac12$O$_2$(g) $\rightarrow$ CO(g)\par
(c) CO(g) + $\dfrac12$ O$_2$(g) $\rightarrow$ CO$_2$(g)\par
Which of the following plots would represent most likely the temperature dependence of $\Delta G^\circ$ ?\par
\begin{qoptions}
\opt{A}{\qfigi{p2_chem_s1_q01_A}}
\opt{B}{\qfigi{p2_chem_s1_q01_B}}
\opt{C}{\qfigi{p2_chem_s1_q01_C}}
\opt{D}{\qfigi{p2_chem_s1_q01_D}}
\end{qoptions}
\qn{2} Potential energy of electron present in 2$^{\text{nd}}$ orbit of Li$^{2+}$ is : ($a_0$ = Radius of 1$^{\text{st}}$ Bohr's orbit) :
\begin{qoptions}
\opt{A}{$-\dfrac{e^2}{4\pi\varepsilon_0a_0}$}
\opt{B}{$-\dfrac{3e^2}{4\pi\varepsilon_0a_0}$}
\opt{C}{$-\dfrac{3e^2}{16\pi\varepsilon_0a_0}$}
\opt{D}{$-\dfrac{9e^2}{16\pi\varepsilon_0a_0}$}
\end{qoptions}
\qn{3} Which of the following is not correct order of reactivity.
\begin{qoptions}
\opt{A}{\qfigi{p2_chem_s1_q03_A}}
\opt{B}{\qfigi{p2_chem_s1_q03_B}}
\opt{C}{\qfigi{p2_chem_s1_q03_C}}
\opt{D}{\qfigi{p2_chem_s1_q03_D}}
\end{qoptions}
\qn{4} Consider the following reaction/s \& choose \textbf{INCORRECT} statement.\par
2NaCl + 2CO$_2$ + 2NH$_3$ + 2H$_2$O $\rightarrow$ (P) + (Q)\par
(P) $\xrightarrow{\Delta}$ (R) + CO$_2$ + H$_2$O
\begin{qoptions}
\opt{A}{In solid salt (P) anion form an infinite chain}
\opt{B}{From compound (Q) NH$_3$ can be obtained by adding Ca(OH)$_2$}
\opt{C}{Compound (P) has symmetrical H-bond}
\opt{D}{MnO$_2$ salt give green coloured bead when heated with compound (R) in presence of air.}
\end{qoptions}
\sect{SECTION-I (ii)}
\qn{5} At 298 K, standard reduction potential is 1.33 V for Cr$_2$O$_7^{2-}$/Cr$^{3+}$.\par
Select the correct statement(s). (If concentration of other species are 1M)\par
Given : $\dfrac{2.303\text{R}\times298}{\text F}=0.06$ V, $E^\circ_{\text{Br}_2/\text{Br}^-}=1.07$ V and $E^\circ_{\text{I}_2/\text{I}^-}=0.54$ V .
\begin{qoptions}
\opt{A}{Cr$_2$O$_7^{2-}$ may oxidise Br$^-$ and I$^-$ at pH = 3}
\opt{B}{Cr$_2$O$_7^{2-}$ may oxidise Br$^-$ and I$^-$ at pH = 0}
\opt{C}{Cr$_2$O$_7^{2-}$ may oxidise I$^-$ but can't oxidise Br$^-$ at pH = 3}
\opt{D}{Cr$_2$O$_7^{2-}$ may oxidise I$^-$ but can't oxidise Br$^-$ at pH = 0}
\end{qoptions}
\qn{6} Choose the \textbf{CORRECT} statement(s) regarding Vanilin\par
\qfig{p2_chem_s2_q06_1}
\begin{qoptions}
\opt{A}{\textbf{[P]} can be obtained as a major product when 3,4-dihydroxy benzaldehyde is treated with one equivalent of NaOH and one equivalent of MeBr.}
\opt{B}{\textbf{[P]} can be obtained as a one of the product when 2-methoxy phenol is treated with chloroform and aq KOH.}
\opt{C}{Positional isomer of \textbf{[P]} can be obtained when 3,4-dihydroxy benzaldehyde is treated with one equivalent of NaOH followed by treatment with one equivalent of MeBr.}
\opt{D}{\textbf{[P]} can be obtained when 3,4-dihydroxy benzaldehyde is treated with two equivalent of NaOH and one equivalent of MeBr followed by protonation.}
\end{qoptions}
\qn{7} Choose correct statement/s about extraction of various metal in metallurgical process.
\begin{qoptions}
\opt{A}{In Mond process the decomposition of nickel carbonyl is accompanied by a large entropy increase (four moles of gas formed from one mole of complex) and hence it takes place spontaneously at higher temperature when T$\Delta$S term becomes significant}
\opt{B}{Crude gold obtained by Mac Arthur Forest cyanide process contain Zn and sometime Pb impurity}
\opt{C}{In Van Arkel-De Boer process the impure element is heated with iodine, producing a volatile iodide.}
\opt{D}{In Kroll process\par TiCl$_4$ + 2Mg $\xrightarrow{\text{under atmosphere of Ar}}$ Ti + 2MgCl$_2$\par Here MgCl$_2$ is leached with H$_2$O}
\end{qoptions}
\sect{SECTION-II (i)}
\paraq{1}{2}
For the reaction, CHCl$_3$(g) + Cl$_2$(g) $\xrightarrow{h\nu}$ CCl$_4$(g) + HCl(g) mechanism is given below\par
Cl$_2$(g) $\underset{k_{-1}}{\overset{k_1}{\rightleftharpoons}}$ 2Cl(g)\par
CHCl$_3$(g) + Cl(g) $\xrightarrow{k_2}$ CCl$_3$(g) + HCl(g)\par
CCl$_3$(g) + HCl(g) $\xrightarrow{k_3}$ CCl$_4$(g)\par
\textbf{Given :}\par
$k_1=4.8\times10^3$ sec$^{-1}$, $k_{-1}=1.2\times10^3$ M$^{-1}$ sec$^{-1}$,\par
$k_2=1.3\times10^{-2}$ M$^{-1}$ sec$^{-1}$, $k_3=2.1\times10^2$ M$^{-1}$ sec$^{-1}$\par
$E_{a_1}$ = 10 kJ/mol, $E_{a_{-1}}$ = 15 kJ/mol,\par
$E_{a_2}$ = 45 kJ/mol, $E_{a_3}$ = 12 kJ/mol
\qn{1} Overall rate constant of reaction, CHCl$_3$(g) + Cl$_2$(g) $\xrightarrow{h\nu}$ CCl$_4$(g) + HCl(g), is \_\_\_\_\_\_ $\times10^{-2}$ M$^{-1/2}$ sec$^{-1}$.
\qn{2} Overall activation energy of reaction, CHCl$_3$(g) + Cl$_2$(g) $\xrightarrow{h\nu}$ CCl$_4$(g) + HCl(g), is \_\_\_\_\_\_\_\_ kJ/mole.
\qfullbegin
\paraq{3}{4}
\qfig{p2_chem_s3_para3_4_F}
\qfullend
\qn{3} Total number of fractions obtained on fractional distillation of all monochloro derivatives product(s) (C)?
\qn{4} Calculate mass of Cl$_2$ required to react with 9.8 gm of compound (B) to form monochloro product. Assuming \% yield = 100\%
\sect{SECTION-II (ii)}
\qn{5} A metal M form sparingly soluble metal sulphate. The logarithm of concentration of dissolved metal ion in a solution saturated in the metal sulphate is plotted below as a function of logarithm of concentration of sulphate ion in solution. The molecular formula of metal sulphate is M$_x$(SO$_4$)$_y$. Then find z $\times$ 10$^{-6}$.\par
$z=\dfrac{x+y}{K_{sp}\text{ of M}_x\text{(SO}_4)_y}$, x and y are minimum integer co-efficient.\par
\qfig{p2_chem_s4_q05_1}
\qn{6} Total number of Pseudohalides ion which are planar and polar.\par
$\overset{\ominus}{\text{C}}$N , $\overset{\ominus}{\text{S}}$CN , $\overset{\ominus}{\text{Se}}$CN , $\overset{\ominus}{\text{O}}$CN , CN$_2^{2-}$ , $\overset{\ominus}{\text{O}}$NC , N$_3^\ominus$ , I$_3^\ominus$ , ICl$_4^\ominus$
\qn{7} Total number of reaction/s in which paramagnetic gaseous product form on heating.\par
(1) KMnO$_4$ $\xrightarrow{\Delta}$\par
(2) MnO$_2$ + conc. HCl $\xrightarrow{\Delta}$\par
(3) MnO$_2$ + conc. H$_2$SO$_4$ $\xrightarrow{\Delta}$\par
(4) KClO$_3$ $\xrightarrow{\Delta}$\par
(5) Pb(NO$_3$)$_2$ $\xrightarrow{\Delta}$\par
(6) CaCO$_3$ $\xrightarrow{\Delta}$\par
(7) NaNO$_3$ $\xrightarrow{\Delta}$\par
(8) NH$_4$NO$_3$ $\xrightarrow{\Delta}$\par
(9) Na$_2$CrO$_4$ + conc. H$_2$SO$_4$ $\rightarrow$
\qn{8} Total number of stereoisomer possible for complex compound \qfigi[0.75]{p2_chem_s4_q08_1}
\qn{9} \qfig{p2_chem_s4_q09_1}
Total number of reagent/s which give redox reaction with solution (A) in acidic medium among the following.\par
Pb(CH$_3$COO)$_2$ , Ba(NO$_3$)$_2$ , NO$_2^-$ , Sn$^{2+}$ , SO$_2$ , Fe$^{2+}$ , H$_2$S , I$^-$ , C$_2$H$_5$OH
\qn{10} How many amino acids are positively charged at pH = 7.\par
(i) Asparagine \hfill (ii) Valine \hfill (iii) Aspartic acid\par
(iv) Glycine \hfill (v) Lysine \hfill (vi) Glutamine\par
(vii) Glutamic acid \hfill (viii) Arginine \hfill (ix) Methionine
\raggedcolumns
\partsec{3}{MATHEMATICS}
\sect{SECTION-I (i)}
\qn{1} Let $\displaystyle\lim_{n\to\infty}\dfrac1n\left\{\left(\sum_{r=1}^n\ell n\left(n^2+r^2\right)\right)-2n\,\ell n\,n\right\}=\ell n2+\dfrac\pi2-2$\par
If\par\resizebox{\linewidth}{!}{$\displaystyle\lim_{n\to\infty}\dfrac1{n^{2m}}\left[\left(n^2+1^2\right)^m\left(n^2+2^2\right)^m\left(n^2+3^2\right)^m\ldots\left(n^2+n^2\right)^m\right]^{\frac1n}=\left(\dfrac{a\sqrt{e\pi}}{e^2}\right)^m$}\par then the value of a is (where m, n $\in$ N)
\begin{qoptions}
\opt{A}{1}
\opt{B}{$\dfrac32$}
\opt{C}{2}
\opt{D}{$\dfrac12$}
\end{qoptions}
\qn{2} The remainder when $\displaystyle\sum_{k=0}^{1012}\left[\dfrac{{}^{1012}C_k}{{}^{2024}C_k}\sum_{r=k}^{2024}{}^rC_k\cdot{}^{2024}C_r\right]$ is divided by 25 is (where ${}^nC_r=\dfrac{n!}{r!(n-r)!}$)
\begin{qoptions}
\opt{A}{15}
\opt{B}{17}
\opt{C}{21}
\opt{D}{11}
\end{qoptions}
\qn{3} A moving particle starts at the point (4, 4) and moves until it hits one of the coordinate axes for the first time. When the particle is at the point (a, b) it moves at random to one of the points (a $-$ 1, b), (a, b $-$ 1) or (a $-$ 1, b $-$ 1) each with probability $\dfrac13$, independent of its previous moves. The probability that it will hit the coordinate axes at (0, 0) is $\dfrac{x}{3^y}$. Where x and y are positive integers such that x is not divisible by 3. Then value of x + y is
\begin{qoptions}
\opt{A}{252}
\opt{B}{253}
\opt{C}{152}
\opt{D}{353}
\end{qoptions}
\qn{4} Let H be the hyperbola having eccentricity `e' and C be the smallest circle touching H and its pair of asymptotes. Ratio of radius of C with the length of semi-conjugate axis of H is.
\begin{qoptions}
\opt{A}{$e+\sqrt{e^2-1}$}
\opt{B}{$e-\sqrt{e^2-1}$}
\opt{C}{$e-\sqrt{e^2+1}$}
\opt{D}{$\dfrac1e$}
\end{qoptions}
\sect{SECTION-I (ii)}
\qn{5} Let A and B be two 3 $\times$ 3 matrices where both A and I + AB are non-singular. If $3A^2B^2+AB=2(AB)^2$ then, which of the following is greater than $|B|$ (where I is identity matrix of order 3 \& $|X|$ denotes determinant of X)
\begin{qoptions}
\opt{A}{1}
\opt{B}{2}
\opt{C}{3}
\opt{D}{4}
\end{qoptions}
\qn{6} Consider f : $\left[0,\dfrac\pi2\right]-\left\{\dfrac\pi4\right\}\rightarrow$ R. Let $f(x)=\displaystyle\lim_{n\to\infty}\prod_{k=0}^n\dfrac{\left((\cot x)^{2^k}+1\right)^2}{\left((\cot x)^{2^{k+1}}+1\right)}$, $\forall x\in\left(0,\dfrac\pi2\right]-\left\{\dfrac\pi4\right\}$ and f$\left(\dfrac\pi4\right)$ = 2 and f(0) = 0, Then which of the following is true?
\begin{qoptions}
\opt{A}{f(x) is increasing in $\left(0,\dfrac\pi4\right)$ \& $\left(\dfrac\pi4,\dfrac\pi2\right)$}
\opt{B}{f(x) is decreasing in $\left(\dfrac\pi4,\dfrac\pi2\right)$}
\opt{C}{f(x) has one local maxima}
\opt{D}{f(x) has three local minima}
\end{qoptions}
\qn{7} If $z=e^{i\pi/13}$ then $\dfrac1{1-z}$ is equal to (Where $i=\sqrt{-1}$)
\begin{qoptions}
\opt{A}{$1-z^2+z^4-z^6+z^8-z^{10}+z^{12}$}
\opt{B}{$1+z^2+z^4+z^6+z^8+z^{10}+z^{12}$}
\opt{C}{$z-z^3+z^5-z^7+z^9-z^{11}$}
\opt{D}{$z+z^3+z^5+z^7+z^9+z^{11}$}
\end{qoptions}
\sect{SECTION-II (i)}
\paraq{1}{2}
Consider a circle C : $x^2+y^2=5$,\par
Plane $\pi_1$ : x $-$ y + 2z + 5 = 0 and\par
Plane $\pi_2$ : 2x + y + z $-$ 5 = 0 and\par
Plane $\pi_3$ : x + 2y $-$ z = 0.\par
Let reflection of curve `C' about plane mirror $\pi_1$ be curve $C_1$. [.] represents G.I.F.
\qn{1} If $\lambda$ is the shortest distance of plane $\pi_2$ from any point on $C_1$ then [$\lambda$] equal to:
\qn{2} Planes $\pi_1$, $\pi_2$, $\pi_3$, form a prism. If `A' is area of cross-section of this prism then [A] equal to
\paraq{3}{4}
Consider $\displaystyle\int_0^{\pi/2}(\sin x)^m(\cos x)^n\,dx$, m, n $\in$ N.\par
We can simplify this integration by substituting sinx = t OR cosx = t.\par
Based on this solve below questions
\qn{3} $I_1=\displaystyle\int_0^1x^6(1-x)^4dx$, $I_2=\displaystyle\int_0^1x^5(1-x)^3dx$\par
then $\dfrac{I_1}{I_2}=\dfrac pq$, where p, q are co-prime positive integers then p + q equal to
\qn{4} Let $I=\displaystyle\int_0^{\pi/4}\dfrac{\left(\cos\left(x+\frac\pi4\right)\right)^{10}\cdot(\sin x)^7}{(\cos x)^{19}}dx=\dfrac{m!\,n!}{2^p\cdot18!}$, m, n, p are natural numbers then smallest value of m + n + p = ?
\sect{SECTION-II (ii)}
\qn{5} Let us consider a complex number $z=x^{3/2}+iy^2$; x $\geqslant$ 0; y $\in$ R such that $2x^5y-x^4y+2x^3y-x^2y+4xy-2y=0$ and $|z|^2=(\text{Re}(z))^{2/3}$. Then number of all distinct value(s) of $|z|$ is
\qn{6} Consider a circle S : $x^2+y^2+x-y=2021$. If (x, y) is called integral points if x and y are both integer. Then number of integral points on S = 0 is
\qn{7} Find number of permutations $x_1,x_2,x_3,x_4,x_5$ of numbers 1, 2, 3, 4, 5 taken all at a time such that the sum of five products $x_1x_2x_3+x_2x_3x_4+x_3x_4x_5+x_4x_5x_1+x_5x_1x_2$ is divisible by 3.
\qn{8} The third derivative of $\dfrac{x^4}{(x-1)(x-2)}$ is $-\dfrac{k}{(x-2)^4}+\dfrac{6}{(x-1)^4}$. The number of $\theta\in[0,2\pi]$ satisfying the equation $\sin^2\theta+2\cos^2\theta=k-96-3\sin\theta\cos\theta$ is \_\_\_\_.
\qn{9} If A = $\begin{bmatrix}6&3\\2&1\end{bmatrix}$ and $10^3<\text{tr}(A^n)<10^4$ where n $\in$ N. The value of n is (where tr(X) denotes trace of matrix X)
\qn{10} Let us consider the curves $9x^2+4y^2=900$, $x^2+y^2-4x-6y-12=0$ and $y^2-6y-6x+51=0$. The maximum number of disjoint regions into which these curve divide xy plane (excluding themselves).
\end{multicols}
\end{document}
